% GENERATED FILE. Edit canonical lessons, then run npm run generate:books.
% canonical-source-hash: fnv1a64:06d890122ed0ecbf
% canonical-lessons: JA-C531-wazato, JA-C531-tsui, JA-C531-muyamini, JA-C531-nozomashii, JA-C531-nantonaku

\chapter{On Purpose or Not}
\label{ch:ja-on-purpose-or-not}
\hlchaptermodality{531}

\begin{chapteropening}
\textbf{By the end of this chapter:} \emph{I can talk about doing things on purpose or by habit.}

Complete the last lesson of chapter 531: I can talk about doing things on purpose or by habit.
\end{chapteropening}

\section[\texorpdfstring{\ja{わざと} (wazato)}{わざと (wazato)}]{\ja{わざと} (wazato) — on purpose; deliberately}

\label{lesson:JA-C531-wazato}



\begin{quote}
Before the new one: say the Japanese for about, then the Japanese for (not) at all.
\end{quote}

\subsection*{You'll want to know: \ja{わざと}}
\textbf{\ja{わざと}} — \emph{wazato} — \textquotedblleft{}on purpose; deliberately\textquotedblright{}.

\begin{quote}
\textbf{\ja{もし} \ja{わざと} \ja{まけたら、みんな} \ja{おこるでしょう。}} — \emph{moshi wazato maketara, minna okoru deshō} — \textquotedblleft{}If you lose on purpose, everyone will be angry.\textquotedblright{}
\end{quote}

\begin{cousinweb}
\textbf{\ja{わざと}} (\emph{wazato}) is a native Japanese word (wago), inherited rather than borrowed from Chinese, and usually written in kana; its older origin is not settled (JMdict).
\end{cousinweb}

\begin{grammarlens}[title={What you've built}]
One more word for saying what happens if.
\end{grammarlens}

\subsection*{Your turn}
\begin{itemize}
\raggedright
  \item \emph{Say it:} \emph{wazato}
  \item \emph{Say it:} \emph{wazato}, once more
  \item \emph{Say it:} \emph{moshi wazato maketara, minna okoru deshō} — If you lose on purpose, everyone will be angry.
  \item \emph{Recall:} say the Japanese for about, then the Japanese for (not) at all, then say \emph{wazato} again
\end{itemize}

\begin{tcolorbox}[breakable,colback=teal!4,colframe=teal!35!black,title={Before you move on}]
What does \ja{わざと} mean? (\textquotedblleft{}On purpose; deliberately\textquotedblright{}.) Say it once more.
\end{tcolorbox}
\section[\texorpdfstring{\ja{つい} (tsui)}{つい (tsui)}]{\ja{つい} (tsui) — without meaning to}

\label{lesson:JA-C531-tsui}



\begin{quote}
Before the new one: say the Japanese for (not) at all, then the Japanese for on purpose.
\end{quote}

\subsection*{You'll want to know: \ja{つい}}
\textbf{\ja{つい}} — \emph{tsui} — \textquotedblleft{}without meaning to\textquotedblright{}.

\begin{quote}
\textbf{\ja{もし} \ja{あまい} \ja{ものが} \ja{あったら、おとうとは} \ja{つい} \ja{たべるでしょう。}} — \emph{moshi amai mono ga attara, otōto wa tsui taberu deshō} — \textquotedblleft{}If there were sweet things about, my younger brother would eat them without meaning to.\textquotedblright{}
\end{quote}

\begin{cousinweb}
\textbf{\ja{つい}} (\emph{tsui}) is usually counted a native Japanese word (wago), inherited rather than borrowed from Chinese: JMdict gives this adverb no kanji spelling, and no dictionary we use records an older form, so its history is not settled (JMdict).
\end{cousinweb}

\begin{grammarlens}[title={What you've built}]
One more word for saying what happens if.
\end{grammarlens}

\subsection*{Your turn}
\begin{itemize}
\raggedright
  \item \emph{Say it:} \emph{tsui}
  \item \emph{Say it:} \emph{tsui}, once more
  \item \emph{Say it:} \emph{moshi amai mono ga attara, otōto wa tsui taberu deshō} — If there were sweet things about, my younger brother would eat them without meaning to.
  \item \emph{Recall:} say the Japanese for (not) at all, then the Japanese for on purpose, then say \emph{tsui} again
\end{itemize}

\begin{tcolorbox}[breakable,colback=teal!4,colframe=teal!35!black,title={Before you move on}]
What does \ja{つい} mean? (\textquotedblleft{}Without meaning to\textquotedblright{}.) Say it once more.
\end{tcolorbox}
\section[\texorpdfstring{\ja{むやみに} (muyamini)}{むやみに (muyamini)}]{\ja{むやみに} (muyamini) — recklessly; without thinking}

\label{lesson:JA-C531-muyamini}



\begin{quote}
Before the new one: say the Japanese for on purpose, then the Japanese for without meaning to.
\end{quote}

\subsection*{You'll want to know: \ja{むやみに}}
\textbf{\ja{むやみに}} — \emph{muyamini} — \textquotedblleft{}recklessly; without thinking\textquotedblright{}.

\begin{quote}
\textbf{\ja{もし} \ja{むやみに} \ja{くすりを} \ja{のんだら、あぶないでしょう。}} — \emph{moshi muyamini kusuri o nondara, abunai deshō} — \textquotedblleft{}If you take medicine without thinking, it will probably be dangerous.\textquotedblright{}
\end{quote}

\begin{cousinweb}
\textbf{\ja{むやみに}} (\emph{muyamini}) is usually counted a native Japanese word (wago), inherited rather than borrowed from Chinese, with the particle \emph{ni}; JMdict marks it as usually written in kana and its kanji spelling as ateji, characters chosen for their sound, and no dictionary we use records an older form, so its history is not settled (JMdict).
\end{cousinweb}

\begin{grammarlens}[title={What you've built}]
One more word for saying what happens if.
\end{grammarlens}

\subsection*{Your turn}
\begin{itemize}
\raggedright
  \item \emph{Say it:} \emph{muyamini}
  \item \emph{Say it:} \emph{muyamini}, once more
  \item \emph{Say it:} \emph{moshi muyamini kusuri o nondara, abunai deshō} — If you take medicine without thinking, it will probably be dangerous.
  \item \emph{Recall:} say the Japanese for on purpose, then the Japanese for without meaning to, then say \emph{muyamini} again
\end{itemize}

\begin{tcolorbox}[breakable,colback=teal!4,colframe=teal!35!black,title={Before you move on}]
What does \ja{むやみに} mean? (\textquotedblleft{}Recklessly; without thinking\textquotedblright{}.) Say it once more.
\end{tcolorbox}
\section[\texorpdfstring{\ja{のぞましい} (nozomashii)}{のぞましい (nozomashii)}]{\ja{のぞましい} (nozomashii) — desirable; best}

\label{lesson:JA-C531-nozomashii}



\begin{quote}
Before the new one: say the Japanese for without meaning to, then the Japanese for recklessly.
\end{quote}

\subsection*{You'll want to know: \ja{のぞましい}}
\textbf{\ja{のぞましい}} — \emph{nozomashii} — \textquotedblleft{}desirable; best\textquotedblright{}.

\begin{quote}
\textbf{\ja{のぞましい} \ja{ほうほうは} \ja{ばあいに} \ja{よります。}} — \emph{nozomashii hōhō wa baai ni yorimasu} — \textquotedblleft{}The best method depends on the case.\textquotedblright{}
\end{quote}

An i-adjective.

\begin{cousinweb}
\textbf{\ja{のぞましい}} (\emph{nozomashii}) is a native Japanese word (wago), inherited rather than borrowed from Chinese, built on the same stem as \emph{nozomu}, \textquotedblleft{}to wish for\textquotedblright{}; its kanji was attached later as a kun'yomi (JMdict, KANJIDIC2).
\end{cousinweb}

\begin{grammarlens}[title={What you've built}]
One more word for saying what happens if.
\end{grammarlens}

\subsection*{Your turn}
\begin{itemize}
\raggedright
  \item \emph{Say it:} \emph{nozomashii}
  \item \emph{Say it:} \emph{nozomashii}, once more
  \item \emph{Say it:} \emph{nozomashii hōhō wa baai ni yorimasu} — The best method depends on the case.
  \item \emph{Recall:} say the Japanese for without meaning to, then the Japanese for recklessly, then say \emph{nozomashii} again
\end{itemize}

\begin{tcolorbox}[breakable,colback=teal!4,colframe=teal!35!black,title={Before you move on}]
What does \ja{のぞましい} mean? (\textquotedblleft{}Desirable; best\textquotedblright{}.) Say it once more.
\end{tcolorbox}
\section[\texorpdfstring{\ja{なんとなく} (nantonaku)}{なんとなく (nantonaku)}]{\ja{なんとなく} (nantonaku) — somehow; vaguely}

\label{lesson:JA-C531-nantonaku}



\begin{quote}
Before the new one: say the Japanese for recklessly, then the Japanese for desirable.
\end{quote}

\subsection*{You'll want to know: \ja{なんとなく}}
\textbf{\ja{なんとなく}} — \emph{nantonaku} — \textquotedblleft{}somehow; vaguely\textquotedblright{}.

\begin{quote}
\textbf{\ja{もし} \ja{なんとなく} \ja{ねつが} \ja{あったら、やすんで} \ja{ください。}} — \emph{moshi nantonaku netsu ga attara, yasunde kudasai} — \textquotedblleft{}If you feel a little feverish, please rest.\textquotedblright{}
\end{quote}

\begin{cousinweb}
\textbf{\ja{なんとなく}} (\emph{nantonaku}) is a native Japanese word (wago), inherited rather than borrowed from Chinese, and usually written in kana; its older origin is not settled (JMdict).
\end{cousinweb}

\begin{grammarlens}[title={What you've built}]
One more word for saying what happens if.
\end{grammarlens}

\subsection*{Your turn}
\begin{itemize}
\raggedright
  \item \emph{Say it:} \emph{nantonaku}
  \item \emph{Say it:} \emph{nantonaku}, once more
  \item \emph{Say it:} \emph{moshi nantonaku netsu ga attara, yasunde kudasai} — If you feel a little feverish, please rest.
  \item \emph{Recall:} say the Japanese for recklessly, then the Japanese for desirable, then say \emph{nantonaku} again
\end{itemize}

\begin{tcolorbox}[breakable,colback=teal!4,colframe=teal!35!black,title={Before you move on}]
What does \ja{なんとなく} mean? (\textquotedblleft{}Somehow; vaguely\textquotedblright{}.) Say it once more.
\end{tcolorbox}
